\cotefile{156-5}
\scanpage{10}\mspage{9}
en question sont dans $\Phi$, i.e. que $\Phi$ est \unsure{fermé|fermée} dans $\Phi\cup\Psi$ (et itou pour le cas $A\in\Psi$, avec $\Psi$ fermé dans $\Phi\cup\Psi$). Donc il \unsure{vient|reste}

\[
\begin{aligned}
2^{\circ})\quad &\deleted{\unsure{\Psi}}\Phi\ \text{fermé dans }\Phi\cup\Psi:\quad A\in\Phi,\ B\in\Psi,\\
&B\subset\deleted{\unsure{B}}A\ \Longrightarrow\ B\in\deleted{B}\Phi.
\end{aligned}
\]
Mais on aura $\deleted{\unsure{\lvert B\rvert}}$ $B\subset\lvert\Phi\rvert\cap\lvert\Psi\rvert$, donc $B\in\Phi\Rightarrow B\deleted{\unsure{\subset}}\in\deleted{\unsure{\Phi}}\Psi$ OK.

\uline{\unsure{Définition|Definition}} On dit alors que $\Phi$ et $\Psi$ \uline{compatibles}

\uline{Corollaire 1} Quand ces conditions sont satisfaites, $\lvert\deleted{A}\Phi\rvert\cap\lvert\Psi\rvert$ est une partie de $\lvert\Phi\rvert$ qui est $\Phi$-saturée, et une partie de $\Psi$ qui est $\Psi$-saturée, et les deux relations d'\unsure{équivalence|équilence} induites par $\lvert\deleted{\unsure{B}}\Phi\rvert$ et par $\lvert\Psi\rvert$ sont les $=\mathrm{s}$.

Mais attention, cette dernière condition ne suffit nullement à \unsure{impliquer} que $\Phi$ et $\Psi$ compatibles. Exemple : $\Phi$ et la famille \inserted{$\Phi^{\circ}$} des $A^{\circ}$ ont $=$ support, et définissent la $=$ relation d'équiv. sur $\lvert\Phi\rvert$, mais elles ne sont compatibles que si elles sont \uline{égales} \deleted{(car} i.e. $\Phi$ formé d'\unsure{ens.} deux à deux disjoints. En effet \deleted{(car la relation d'ordre sur $\Phi$ et sur $\Phi^{\circ}$ n'est la $=$ que si}

\uline{Corollaire 2} Supposons $\Phi$ et $\Psi$ compatibles. Alors $\lvert\Phi\rvert\subset\lvert\Psi\rvert$ \unsure{sss|ssi} $\Phi\subset\Psi$, $\lvert\Phi\rvert=\lvert\Psi\rvert$ sss $\Phi=\Psi$.
